\section{An elementary dual-number derived classification}
\label{app:dual-number-classification}

This appendix proves the classification used in
Theorem~\ref{thm:rank-formula}, retaining the finiteness conditions needed
for that application. The classification is established mathematics;
see Keller--Yang--Zhou~\cite[Example 3.7(b)]{keller}. The argument here
explains why the possible unbounded projective resolution introduces
only simple-module summands, rather than additional unknown types.

Let $A=\F[\epsilon]/(\epsilon^2)$. A cochain complex has differential
raising its degree by one. A finite free interval $J_l$ has $l$ free
rank-one terms and consecutive differentials $\epsilon$. A
left-infinite interval ending in degree $b$ is
\[
 \cdots\xrightarrow{\epsilon}A
 \xrightarrow{\epsilon}A
 \xrightarrow{\epsilon}\underset{b}{A}\longrightarrow0.
\]
It is exact below $b$ and has cohomology $A/(\epsilon)=\F$ in degree
$b$. It therefore represents a shift of the simple module.

\begin{proposition}[Classification needed for continuation ranks]
\label{prop:derived-classification}
Every bounded complex $D$ of finite-dimensional $A$-modules is
quasi-isomorphic to a finite direct sum of homological shifts of finite
free intervals $J_l$ and of the simple module $\F$.
\end{proposition}

\begin{proof}
We give the construction in three steps.

\paragraph{Step 1: a degreewise finite free resolution.}
There is a bounded-above complex $P$ of finite free $A$-modules and a
quasi-isomorphism $f:P\to D$. The following downward construction
establishes the precise finiteness assertion. Suppose $P$ and $f$ have
been constructed in degrees greater than $i$. The degree-$i$ term of
the partially built mapping cone is $D^i\oplus P^{i+1}$, whose cycles
are the finite $A$-module
\[
 Z_i=\bigl\{(d,p):
 d_Dd+f^{i+1}p=0,\quad d_P^{i+1}p=0\bigr\}.
\]
Choose a finite free module $P^i$ surjecting onto $Z_i$, and define
$f^i$ and $d_P^i$ to be its two component maps. The defining equations
of $Z_i$ give the chain-map identity and $d_P^{i+1}d_P^i=0$.
The new boundaries from $P^i$ in the cone cover all its degree-$i$
cycles. Begin above the highest nonzero degree of $D$ and continue
downward. Each cone degree becomes exact, with its outgoing maps
unaffected by later steps. Hence the completed cone is acyclic.
The resulting $P$ is bounded above and has finite free rank in every
degree, as claimed.

\paragraph{Step 2: minimalization and the exact left tail.}
Over the local ring $A$, cancel every unit entry of the differential,
working downward from the highest degrees. A unit pivot splits off a
contractible pair $A\xrightarrow{\sim}A$. Each pair of adjacent degrees
has finite rank, so only finitely many cancellations are needed there.
Once a differential has all its entries in the radical $(\epsilon)$,
later changes of basis cannot make their residues nonzero. The
descending procedure is locally finite: a fixed degree and its adjacent
matrices stabilize after finitely many operations. It therefore gives
a bounded-above minimal free complex $P'$ quasi-isomorphic to $P$.
This also follows by tracking the elementary cancellation equivalences
degree by degree; no infinite sum of maps in a fixed degree is required.

Write $P'^i=A^{r_i}$. Every differential is now uniquely of the form
\[
 d_{P'}^i=\epsilon B_i,
 \qquad B_i:\F^{r_i}\longrightarrow\F^{r_{i+1}}.
\]
Multiplication by $\epsilon B_i$ has binary rank $\rank(B_i)$:
it maps the constant part to the $\epsilon$ part and kills the latter.
Consequently
\begin{equation}\label{eq:minimal-tail-rank}
 \dim_{\F}H^i(P')
  =2r_i-\rank(B_{i-1})-\rank(B_i).
\end{equation}
Because $D$ is bounded, the left side is zero for every sufficiently
negative $i$. Each rank on the right is at most $r_i$, so exactness in
degree $i$ forces
\[
 \rank(B_{i-1})=\rank(B_i)=r_i.
\]
At two consecutive exact degrees this gives $r_i=r_{i+1}$ and makes
$B_i$ invertible. Hence the sufficiently far-left tail has one constant
finite rank $r$, with every $B_i$ invertible. Choose bases recursively
to make all these tail matrices identities. The case $r=0$ simply
means that the minimal complex is bounded on both sides.

\paragraph{Step 3: one finite interval decomposition.}
Choose a cut degree $a$ within that identity tail. Only finitely many
nonzero degrees lie at or to the right of $a$, and their dimensions
are finite. Apply the elementary finite interval decomposition from
Theorem~\ref{thm:interval} to the sequence of binary matrices from
degree $a$ to the highest degree of $P'$.

Every interval born strictly after $a$ remains a finite interval.
Every interval born at $a$ extends to the left along one of the
identity tail coordinates. To see this explicitly, extend the
chosen change of basis in degree $a$ backward through the identity
tail; all its matrices remain identities. Thus these intervals
become left-infinite intervals, with no additional couplings.
There are finitely many intervals altogether because the finite
sequence being decomposed has finite total dimension.

Reinsert the factor $\epsilon$ in the differentials. This expresses
$P'$ as a finite direct sum of finite free $\epsilon$-intervals and
left-infinite $\epsilon$-intervals. Each left-infinite interval is
quasi-isomorphic to a shifted simple module, as observed above.
Composing with the quasi-isomorphisms in the first two steps proves
the proposition.
\end{proof}

Only existence is needed in the continuation argument. The recognizer
does not construct this resolution for a suffix or attempt to process
an infinite complex. It computes the finite barcode of an already
stored pure component, then applies the rank identity which the
classification proves for every possible suffix.

The distinction between free and arbitrary module continuations matters.
If $D$ is perfect, only finite free intervals occur, so the coefficient
$a$ in~\eqref{eq:continuation-rank-formula} vanishes. In an arbitrary
bounded module complex, shifted simple modules can occur. The geometric
cutoff-two theorem did not assume their absence: it proved that their
number $a$ is even, using the unreduced circle-space parity of the
actual closure. This is why that theorem applies to every square-zero
dot polynomial, including ones whose action is not free on a resolution.
